\documentclass[10pt, a4paper]{article}
\usepackage{preamble}

\title{Complex Analysis II}
\author{Luke Phillips}
\date{January 2026}

\begin{document}

\maketitle

\newpage

\tableofcontents

\newpage

\section{The Complex Plane and Complex Functions}

\subsection{Complex Numbers}

Complex numbers are of the form $z = x + iy$ where $i$ is the imaginary unit,
$i = \sqrt{-1}$.
The set of complex numbers is denoted $\C$.
We use an Argand diagram to visualise complex numbers $x + iy$ with $(x, y)$.

For $z = x + iy$ we have the real part of $z$,
$\mathrm{Re}(z) = x$,
and the imaginary part of $z$,
$\Im(z) = x$.

Addition of complex numbers,
$z_1 = x_1 + iy_1$,
$z_2 = x_2 + iy_2$
\[
z_1 \pm z_2 = (x_1 \pm x_2) + i(y_1 \pm y_2)
\]
the real parts are added together and the imaginary parts are added together.

Multiplication of complex numbers,
$z_1 = x_1 + iy_1$,
$z_2 = x_2 + iy_2$
\[
z_1z_2 = (x_1 + iy_1)(x_2 + iy_2) = (x_1x_2 - y_1y_2) + i(x_1y_2 + x_2y_1)
\]
just expand and use $i ^ 2 = -1$.

Looking at
\[
(x + iy)(x - iy) = x ^ 2 + ixy - ixy - i ^ 2y ^ 2 = x ^ 2 + y ^ 2 \in \R.
\]
We call the conjugate of a complex number as follows,
given $z = x + iy$ the conjugate is
\[
\bar{z} = x - iy = \mathrm{Re}(z) - i\mathrm{Im}(z).
\]
We can see that
\[
z\bar{z} = \mathrm{Re}(z) ^ 2 + \Im(z) ^ 2 \in \R
\]
which leads us to
\[
|z| = \sqrt{z\bar{z}} = \sqrt{\mathrm{Re}(z) ^ 2 + \Im(z) ^ 2}
\]
the modulus
(or absolute value)
of $z$.

Using the conjugate we can see
\[
\frac{z_1}{z_2} = \frac{x_1 + iy_1}{x_2 + iy_2} = \frac{(x_1 + iy_1)(x_2 - iy_2)}{(x_2 + iy_2)(x_2 - iy_2)} = \frac{x_1x_2 + y_1y_2}{x_2 ^ 2 + y_2 ^ 2} + i\frac{x_2y_1 - x_1y_2}{x_2 ^ 2 + y_2 ^ 2}.
\]

\begin{example}
    Find $\frac{3}{1 + i} - (3 + 2i)$.

    \begin{solution}
        \[
        \frac{3}{1 + i} - (3 + 2i) = \frac{3(1 - i)}{(1 + i)(1 - i)} - (3 + 2i) = \frac{3 - 3i}{2} - 3 - 2i = -\frac{3}{2} - \frac{7}{2}i.
        \]
    \end{solution}
\end{example}

Conjugation is nothing but reflection with respect to the real axis.

\begin{lemma}[Important Properties of Complex Numbers]
    \begin{enumerate}[label = (\roman*)]
        \item $z_1z_2 = 0 \iff z_1 = 0\text{ or } z_2 = 0$.

        \item $|z| = \sqrt{z\bar{z}}$.

        \item $\overline{(zw)} = \bar{z}\bar{w}$.

        \item $\Re(z) = \frac{z + \bar{z}}{2}$ and $\Im(z) = \frac{z - \bar{z}}{2i}$.

        \item $z ^ {-1} = \frac{\bar{z}}{|z| ^ 2}$.
    \end{enumerate}
\end{lemma}

Using polar coordinates and the Argand diagram we can write complex numbers in polar coordinates.
We define for any $z \in \C$
\[
r \coloneqq \sqrt{x ^ 2 + y ^ 2} = \sqrt{\Re(z) ^ 2+ \Im(z) ^ 2} = |z|
\]
and any $z \neq 0$
\[
\theta \coloneqq \begin{cases}
    \arctan(y / x) & x > 0, \\
    \pi / 2 & x = 0, y > 0, \\
    -\pi / 2 & x = 0, y < 0, \\
    \arctan(y / x) + \pi &x < 0, y > 0, \\
    \arctan(y / x) - \pi &x < 0, y < 0.
\end{cases}
\]
$\theta$ is the anticlockwise angle measured from the real axis,
$\theta$ is called the argument of $z$,
denoted $\arg(z)$.

We have the following polar coordinates for $z$:
\[
z = x + iy = r\cos(\theta) + ir\sin(\theta) = r(\cos(\theta) + i\sin(\theta)).
\]

\begin{example}
    $z = i$.

    \[
    r = |z| = \sqrt{0 ^ 2 + 1 ^ 2} = 1
    \]
    \[
    \theta = \pi / 2.
    \]
    \[
    i = 1(\cos(\pi / 2) + i\sin(\pi / 2)) = \cos(\pi / 2) + i\sin(\pi / 2).
    \]

    $z = 1 + i$
    \[
    |z| = \sqrt{1 ^ 2 + 1 ^ 2} = \sqrt{2},
    \]
    $\theta = \frac{\pi}{4}$,
    \[
    z = \sqrt{2}(\cos(\pi / 4) + i\sin(\pi / 4)).
    \]
\end{example}

\textbf{Important:}
$\arg(z)$ is only defined up to multiples of $2\pi$.

$\arg(i) = \frac{\pi}{2} + 2\pi k$,
$k \in \Z$,
not a function.

To define an argument we need a range of angles.

The principal value of $\arg(z)$ is the value in the interval $(-\pi, \pi]$ denoted $\Arg(z)$.

For example:
$\Arg(i) = \frac{\pi}{2}$,
$\Arg(-1) = \pi$.

\begin{lemma}[Properties of argument]
    We have the following properties of the argument:
    
    \begin{enumerate}[label = (\roman*)]
        \item $\arg(z_1z_2) = \arg(z_1) + \arg(z_2)\mod{2\pi}$.

        \item $\arg(1 / z) = -\arg(z)\mod{2\pi}$.

        \item $\arg(\bar{z}) = -\arg(z)\mod{2\pi}$.
    \end{enumerate}
\end{lemma}
\textit{Note:
the proof relies on why we went to polar coordinates:
multiplication is clear.}

\begin{lemma}
    Geometrically,
    multiplication in $\C$ is given by a dilated rotation;
    i.e.,
    if $z_1 = r_1(\cos(\theta_1) + i\sin(\theta_1))$ and $z_2 = r_2(\cos(\theta_2) + i\sin(\theta_2))$ then
    \[
    z_1z_2 = r_1r_2(\cos(\theta_1 + \theta_2) + i\sin(\theta_1 + \theta_2)).
    \]
    In particular,
    multiplying by $z_2$ constitutes an
    (anticlockwise)
    rotation of $z_1$ by $\theta_2$ degrees,
    followed by a dilation with factor $r_2$.

    \begin{proof}
        \begin{align*}
            z_1z_2 &= r_1(cos(\theta_1) + i\sin(\theta_2))r_2(\cos(\theta_2) + i\sin(\theta_2)) \\
            &= (r_1\cos(\theta_1) + ir_1\sin(\theta_1))(r_2\cos(\theta_2) + ir_2\sin(\theta_2)) \\
            &= (r_1r_2\cos(\theta_1)\cos(\theta_2) - r_1r_2\sin(\theta_1)\sin(\theta_2)) \\
            &+ i(r_1r_2\sin(\theta_1)\cos(\theta_2) + r_1r_2\cos(\theta_1)\sin(\theta_2)) \\
            &= r_1r_2[(\cos(\theta_1)\cos(\theta_2) - \sin(\theta_1)\sin(\theta_2)) + i(\sin(\theta_1)\cos(\theta_2) + \cos(\theta_1)\sin(\theta_2))] \\
            &= r_1r_2(\cos(\theta_1 + \theta_2) + i\sin(\theta_1 + \theta_2)).
        \end{align*}
    \end{proof}
\end{lemma}
\textit{Note:
this proves $\arg(z_1z_2) = \arg(z_1) + \arg(z_2)\mod{2\pi}$,
but also:
$\frac{1}{z} = \frac{\bar{z}}{|z| ^ 2}$ so $\arg(1 / z) = \arg(\bar{z})\mod{2\pi}$,
or $0 = \arg(z \cdot 1 / z) = \arg(z) + \arg(1 / z)\mod{2\pi}$.}

Useful notation:
write $\cos(\theta) + i\sin(\theta) \coloneqq e ^ {i\theta}$.

Motivated by
\[
e ^ {i\theta_1}e ^ {i\theta_2} = (\cos(\theta_1) + i\sin(\theta_1))(\cos(\theta_2) + i\sin(\theta_2)) = \cos(\theta_1 + \theta_2) + i\sin(\theta_1 + \theta_2) = e ^ {i(\theta_1 + \theta_2)}
\]

\begin{corollary}\phantom{}
    \begin{enumerate}[label = (\roman*)]
        \item $|z_1z_2| = |z_1||z_2|$.
        
        \item De Moivre's formula:
        $(\cos(\theta) + i\sin(\theta)) ^ n = \cos(n\theta) + i\sin(n\theta)$.
    \end{enumerate}

    \begin{proof}
        \begin{enumerate}[label = (\roman*)]
            \item Exercise.

            \item Done with induction.
        \end{enumerate}
    \end{proof}
\end{corollary}

The modulus has the following properties.
\begin{enumerate}[label = (\roman*)]
    \item
    (Triangle inequality)
    $|z_1 + z_2| \leq |z_1| + |z_2|$.

    \item $|z| \geq 0$ and $|z| = 0 \iff z = 0$.

    \item $\max(|\Re(z)|, |\Im(z)|) \leq |z| \leq |\Re(z)| + |\Im(z)|$.
\end{enumerate}

We can use $\arg(z)$,
$|z|$ to describe geometric sets.
$\mathbb{D} \coloneqq \{z \in \C\,:\, |z| < 1\}$ -
Unit disc.

$\mathbb{H} = \{z \in \C\,:\, \Im(z) > 0\}$ -
Upper half plane
(without the $x$ axis).

$\mathbb{H}_{R} \coloneqq \{z \in \C,\:\, \Re(z) > 0\}$ -
Right half plane
(without the $y$ axis).

$\mathbb{H}_{L} \coloneqq \{z \in \C,\:\, \Re(z) < 0\}$ -
Left half plane
(without the $y$ axis)

$\arg(z) = \frac{\pi}{4}$ -
Ray of angle $\pi / 4$.

$\{z \in \C\,:\, |z - i| = 4\}$ -
Circle centred at $i$ with radius $4$.

\begin{example}
    What is the set $|z - i| = |z + i|$?

    \begin{solution}
        The real line.

        $z = x + iy$
        \begin{align*}
            |x + i(y - 1)| = |x + i(y + 1)| &\iff |x + i(y - 1)| ^ 2 = |x + i(y + 1)| ^ 2 \\
            &\iff x ^ 2 + (y - 1) ^ 2 = x ^ 2 + (y + 1) ^ 2 \\
            &\iff 4y = 0 \implies y = 0.
        \end{align*}

        With $z$:
        one can show
        \[
        |z - \omega| ^ 2 = |z| ^ 2 - 2\Re(z\bar{\omega}) + |\omega| ^ 2.
        \]
    \end{solution}
\end{example}

\subsection{Complex Functions}

Complex functions are $f : \C \to \C$.
To visualise the graph of complex functions we need four dimensions.
We can always write
\[
f(z) = \Re(f(z)) + i\Im(f(z)).
\]
This is not enough to visualise it.

To understand complex functions we usually try to see how they take sets from $\C$ to sets in $\C$.
We will explore
"elementary"
functions like polynomials,
exponentials,
trigonometric functions,
logarithms and powers.

\subsection{The function \texorpdfstring{$f(z) = z ^ n$}{} with \texorpdfstring{$n \in \N$}{}}
To understand $z ^ n$ we will use polar representation:
if $z = re ^ {i\theta}$.
By our lemma,
\[
z ^ n = r ^ ne ^ {in\theta}
\]
if we fix $r$,
i.e. we are on a circle $z ^ n$ dilates our radius by power of $n$
($r \mapsto r ^ n$) and multiplies our argument by a factor of $n$
($\theta \to n\theta\mod{2\pi}$).

A segment of angular length $\eta$ on a circle of radius $r$ becomes an angular segment of length $n\eta\mod{2\pi}$ if $n\eta < 2\pi$ and the full circle if $n\eta \geq 2\pi$ with radius $r ^ n$.

If $n\eta > 2\pi$,
$z ^ n$ will not be injective.

We can slice the world to wedges of size $2\pi / n$ on which $z ^ n$ is injective.
For instance
\[
S_{k}^{(n)} = \{z \in \C\,:\,-\pi < 2\pi k / n < \Arg(z) \leq -\pi + 2\pi(k + 1) / n\}
\]
where $k = 0, \dotsc, n - 1$.
Or
\[
R_{k}^{(n)} = \{z \in \C\,:\, 2\pi k / n\mod{2\pi} < \Arg(z) \leq 2\pi(k + 1) / n\mod{2\pi}\}
\]
$k = 0, \dotsc, n - 1$.







\end{document}